\section{从受禅到五丈原：同一年份里的三种国家}

这条时间线把魏、蜀汉、孙吴和它们面对的边缘放在同一页上。它不把册命、
行军、继承或使节当成彼此独立的新闻：一地能够调用的粮、道路、官署和
盟友，正是另一地必须计算的限制。账本中同一事件簇的准备、执行和复核记录
仍保留在附录索引；这里叙述其可确认的主节点。

\subsection{220--223：名号落在旧有的军政网络上}

二二〇年正月，曹操去世、曹丕继位。
曹丕首先要证明的不是自己有一套崭新的国家，而是魏王国控制的军队、朝廷和
北方资源不会因继承而断裂。十月的
汉魏禅代遂把这一事实改写成典礼；
献帝退位而曹丕建魏，\eventanchor{EH-0220-HAN-ABDICATION}{献帝禅位于曹丕}。诏册能够确定名义交接的顺序，却不能自行说明起草过程、
群臣的余地或所谓“自愿”的程度。

随后的魏定都洛阳把尚书、禁军和仓储
重新系到旧都。洛阳的价值不在于一座空城的象征，而在于它把中原道路、任命和
漕运可能汇合起来；也正因如此，魏此后向淮南、关中和辽东的命令都要从这一
中枢穿出。

西面并没有因此接受魏的时间表。二二一年，
刘备在成都称帝，以汉的继承者
立名；同时他集结益州军东下，把
荆州与关羽之死造成的损失放到对吴作战的议程上。名义和兵力在这里互相需要：
成都的称号不能替代长江沿线的兵站，东征也不能替代盆地的稳定。

二二二年，曹丕先以册封孙权为
吴王、索取质子来试探江东能否被纳入魏的秩序；江面上的力量却不受册文支配。
夷陵、猇亭之战中，陆逊击破东下的蜀军，刘备
退向白帝。战果可以确认，具体火攻过程、兵数和每一次指挥则见于立场不同的
叙事，不能把后出的戏剧性细节当作战场记录。夷陵改变的首先是选择：蜀汉不能
同时把主要兵力押在江东和北方，孙吴也不愿把魏的册封当成安全。

刘备死于永安，刘禅继位，诸葛亮受遗诏
辅政。蜀汉要先处理继承和后方，不能以复仇延续东征。邓芝的出使使
蜀吴复通；这不是感情上的和解，
而是两国把抗魏的压力分担到盟约上的选择。次年曹丕亲至广陵，
因江冰、舟师和渡江条件撤回，也让
魏看见册命之外的河流和季节。

\subsection{225--229：把后方、山口与江面做成国家能力}

诸葛亮在二二五年的南征，平定益州郡、
永昌一带的叛乱并重置地方关系。南中并非成都可以随意支取的仓库；现存材料能
说明军事和地方重整的结果，却不足以把不同地区的服从写成一次完成。其直接
后果是蜀汉可把更多注意力转回汉中，代价则是国家能力仍须经地方中介来实现。

北方也发生继承。二二六年曹丕去世，曹叡
继位。蜀汉在二二七年由诸葛亮
驻汉中并上《出师表》来组织北伐准备；这份文本照亮的是动员的政治语言，
不是一张足以量出全军粮秣的预算表。

二二八年使三条走廊同时紧张。诸葛亮
出祁山，南安、天水、安定
一度响应，魏调援军西上；街亭失守后，
蜀军撤回汉中并处置失军者。山口的失误并不只
是将领品评：一支从盆地推出的军队必须保持粮道、据点和地方响应，任一环节
收缩都会改变前线。

同年魏以司马懿急行军围新城，
孟达之乱失败；孙吴方面，
陆逊在石亭击败曹休。洛阳可以在近畿和淮南
调度力量，却也因此必须在西线、南线之间分配注意。二二九年，孙权
在武昌称帝，又
迁都建业；独立国号和下游都城共同表明，
吴要把长江水路而非魏的册命作为中枢的支点。蜀汉则
取得武都、阴平，为汉中西北增加一层屏障，
但没有消除向陇右输送的距离。

\subsection{230--234：道路拒绝服从，联盟也有边界}

魏的反攻在二三〇至二三一年暴露了同一种约束：
曹真等多路准备攻蜀，山道因大雨受阻。
地理不是静止背景；栈道、天气和转输会迫使一项已经作出的战略决定停下来。吴
在二三〇年派卫温、诸葛直出海，
求夷洲、亶洲而多病死。史书所记的
远航和损失可用来追问海上动员，不能据此断定航程、岛屿对应或沿海社会已经
被长久控制。

二三一年诸葛亮再出祁山，因粮运
不继撤军；负责转输的李严随即
被弹劾废黜。这两件事相连，却不宜
简化成一个人的过失：成都到前线的道路把军事判断和行政责任绑在一起。孙吴
想借辽东牵制魏，先由公孙渊
通吴受封燕王，次年却
杀吴使、转向魏。远方册封的脆弱，
正与江东对长江的控制不同。

二三四年是一次并行的收束。孙权
攻合肥新城而退；诸葛亮
驻五丈原、病逝，蜀军撤回，其后
蒋琬录尚书事。这些结果并不意味着
三个国家的战争停止，而是把问题交给新的辅政者和新的继承安排：谁能维持一条
道路，谁又能让中枢的命令在道路尽头仍被执行。

\begin{sourcenote}[这一段共同年代的材料边界]
年序主要依《三国志》魏、蜀、吴书及《资治通鉴》；它们能把同年发生于不同
政权的事并列，却各有编纂框架。南中、夷洲和辽东尤其需要《华阳国志》、裴注
所引材料与区域研究互相校读；远征结果不等于对地方或海域的长期控制。
\end{sourcenote}
